%% Notes for lesson 01 — example of how to use the notes file
\begin{campo}{1. What are some questions that he Bible answers?}
    The book raises a few questions, such as:
    \begin{itemize}
        \item How did life begin? 
        \item Why are we here? 
        \item Why do innocent people suffer? 
        \item What happens when a person dies? 
        \item Whay is there so much war? 
        \item What will happen to the earth in the future?
    \end{itemize}
    I hope to answer all of these questions by the time I complete this book.
\end{campo}

\begin{campo}{2. How can the Bible help us to enjoy our daily life?}
    The Bible has several teachings about ordinary thing in our lives, such as how to deal with stress.
\end{campo}

\begin{campo}{Text in focus: 2 Timothy 3:16-17}
The lesson cites \rb{2Ti 3:16}. The context (\rb{3:14-17}) is Paul's counsel to
Timothy, who had known the \enquote{holy writings} from infancy (\rb{3:15}).

\interlinear[modo=leitura, titulo={2Ti 3:16--17 · SBLGNT}, legenda=false]{dados/biblia/2tm03_16-17.grc.tsv}

\begin{lexico}{θεόπνευστος}{G2315}{1}
An adjective built from \gr{θεός} (God) and the verb \gr{πνέω} (to breathe):
\enquote{God-breathed}. It occurs only here in the greek scriptures. The verse lists four
purposes with \gr{πρός}: teaching, reproof, setting things straight, and
discipline in righteousness. \gr{ἐπανόρθωσις} and \gr{ἄρτιος} (complete,
fitted) likewise occur only once in the greek scriptures.
\end{lexico}

\end{campo}

\begin{campo}{3. The Bible can guide us}
    \begin{itemize}
        \item He views the Bible as a guide for illuminating the path he ought to pursue.
        \item The Bible is the word of Jehovah God; it contains messages of everlasting life. 
    \end{itemize}
\end{campo}
\begin{campo}{4. The Bible can answer our questions}
    \begin{itemize}
        \item Why do disabilities occur in children, and why does God allow such suffering to persist? 
        \item Studying the Bible provides her with all the right answers.
        \item How does one accomplish God's will, and how might one identify the true religion?
    \end{itemize}
\end{campo}

\begin{campo}{5. You can enjoy reading the Bible}
    \begin{itemize}
        \item Generally, they do not enjoy rading books.
        \item This is beacause the Bible is much more than a mere book.
        \item Yes, I'm interested in this.
    \end{itemize}
\end{campo}

\begin{campo}{}
\end{campo}

\begin{campo}{Central idea}
(Write the answer of the lesson here, in your own words.)
\end{campo}

\EBlivcampo[3]{To look into}{}
\EBlivcampo[3]{Application}{}
